\section*{Chapitre \ref{ch:g10:inertia} --- Forces et le principe d'inertie}

\begin{solution}{exo:g10:inertia:1}
(a) \omterm{def:g10:universal-gravitation:weight}{Poids} (à distance), \omterm{def:g10:inertia:inventory}{réaction normale} de la table
(contact). (b) \omterm{def:g10:universal-gravitation:weight}{Poids} (à distance), \omterm{def:g10:inertia:inventory}{tension} du fil
(contact). (c) \omterm{def:g10:universal-gravitation:weight}{Poids} (à distance), \omterm{def:g10:inertia:inventory}{réaction normale}
de la glace (contact)~; \omterm{def:g10:inertia:inventory}{frottement} négligeable.
(d) \omterm{def:g10:universal-gravitation:weight}{Poids} seulement --- rien ne touche la pomme.
\end{solution}

\begin{solution}{exo:g10:inertia:2}
Livre~: $P = 0.250 \times 9.81 \approx \qty{2.45}{N}$, une flèche
d'environ \qty{2.5}{cm} de long. Personne~:
$P = 60 \times 9.81 \approx \qty{589}{N}$ --- une flèche de
\qty{5.89}{m}. Choisir une échelle pour faire tenir la plus grande
\omterm{def:g10:inertia:force}{force} sur la page.
\end{solution}

\begin{solution}{exo:g10:inertia:3}
\omterm{def:g10:inertia:force}{Point d'application} $A$~; direction horizontale, pointant vers
la gauche~; intensité $3.5 \times 2 = \qty{7.0}{N}$.
\end{solution}

\begin{solution}{exo:g10:inertia:4}
\omterm{def:g10:universal-gravitation:weight}{Poids} $\vect P$~: vertical, vers le bas,
$P = 1.2 \times 9.81 \approx \qty{11.8}{N}$. \omterm{def:g10:inertia:inventory}{Réaction
normale} $\vect R$ de l'étagère~: même droite, vers le haut,
$R = \qty{11.8}{N}$ (le dictionnaire est au repos, donc les deux se
\omterm{def:g10:inertia:compensate}{compensent}).
\end{solution}

\begin{solution}{exo:g10:inertia:5}
(a) Faux~: le \omterm{def:g10:inertia:inventory}{frottement} arrête les choses~; sans aucune
\omterm{def:g10:inertia:force}{force} le corps garde son \omterm{def:g10:inertia:uniform}{mouvement rectiligne
uniforme}. (b) Vrai (principe d'inertie). (c) Faux~: il est soumis à
des \omterm{def:g10:inertia:force}{forces} qui se \emph{\omterm{def:g10:inertia:compensate}{compensent}}
(p.ex.\ \omterm{def:g10:universal-gravitation:weight}{poids} et réaction). (d) Vrai --- la lecture
contraposée du principe.
\end{solution}

\begin{solution}{exo:g10:inertia:6}
\emph{1.} $v = 24/3.0 = \qty{8.0}{m/s}$. \omterm{def:g10:inertia:force}{Forces}~:
\omterm{def:g10:universal-gravitation:weight}{poids} et \omterm{def:g10:inertia:inventory}{réaction normale}~; elles se
\omterm{def:g10:inertia:compensate}{compensent} (\omterm{def:g10:inertia:uniform}{mouvement rectiligne
uniforme}, \omterm{def:g10:inertia:inventory}{frottement} négligeable). \emph{2.} Le
\omterm{def:g10:inertia:inventory}{frottement} du béton n'est plus négligeable~: l'inventaire
gagne une \omterm{def:g10:inertia:force}{force} vers l'arrière non compensée, et en effet la
vitesse change --- le palet ralentit et s'arrête.
\end{solution}

\begin{solution}{exo:g10:inertia:7}
\omterm{def:g10:universal-gravitation:weight}{Poids} $P = 70 \times 9.81 \approx \qty{687}{N}$ vers le bas~;
\omterm{def:g10:inertia:force}{force} du plancher vers le haut. \omterm{def:g10:inertia:uniform}{Mouvement
rectiligne uniforme}, donc elles se \omterm{def:g10:inertia:compensate}{compensent}~: le
plancher pousse avec \qty{687}{N}. Rien ne change au repos --- le repos
et le mouvement uniforme donnent le même \omterm{def:g10:inertia:diagram}{diagramme des
forces}.
\end{solution}

\begin{solution}{exo:g10:inertia:8}
\emph{1.} Au repos, la \omterm{def:g10:inertia:inventory}{tension} \omterm{def:g10:inertia:compensate}{compense} le
\omterm{def:g10:universal-gravitation:weight}{poids}~: $T = 1.8 \times 9.81 \approx \qty{17.7}{N}$.
\emph{2.} Chaque fil vertical porte la moitié~: $T = \qty{8.8}{N}$.
\end{solution}

\begin{solution}{exo:g10:inertia:9}
Freinage~: le corps du passager garde sa vitesse (inertie) tandis que
le bus perd la sienne, donc le passager se déplace en avant
\emph{par rapport au bus}. Virage~: le corps tend à continuer droit
tandis que le bus tourne en dessous, donc il dérive vers l'extérieur
du virage. Dans les deux cas la «~\omterm{def:g10:inertia:force}{force}~» est fictive~; ce
qui manque est une vraie \omterm{def:g10:inertia:force}{force} (adhérence, main courante)
pour changer la vitesse du passager avec celle du bus.
\end{solution}

\begin{solution}{exo:g10:inertia:10}
\omterm{def:g10:inertia:uniform}{Mouvement rectiligne uniforme}, donc \omterm{def:g10:universal-gravitation:weight}{poids} et
résistance de l'air se \omterm{def:g10:inertia:compensate}{compensent}~: traînée
$= P = 85 \times 9.81 \approx \qty{834}{N}$, verticale et vers le
haut.
\end{solution}

\begin{solution}{exo:g10:inertia:11}
\emph{1.} \omterm{def:g10:inertia:uniform}{Mouvement rectiligne uniforme}~: les
\omterm{def:g10:inertia:force}{forces} horizontales se \omterm{def:g10:inertia:compensate}{compensent},
donc la traînée est \qty{300}{N}, opposée à la corde. \emph{2.} Le
\omterm{def:g10:universal-gravitation:weight}{poids} du skieur et la poussée vers le haut de l'eau sur
les skis.
\end{solution}

\begin{solution}{exo:g10:inertia:12}
\emph{1.} $P = 2.4 \times 9.81 \approx \qty{23.5}{N}$~; les
composantes verticales portent le \omterm{def:g10:universal-gravitation:weight}{poids}~:
$2T\cos 40^\circ = P$, donc
$T = 23.5/(2 \times 0.766) \approx \qty{15.4}{N}$. \emph{2.}
$T = P/(2\cos\theta)$ croît sans borne lorsque $\theta \to 90^\circ$
puisque $\cos\theta \to 0$~; des fils exactement horizontaux n'ont
\emph{aucune} composante verticale, donc rien ne porterait le
\omterm{def:g10:universal-gravitation:weight}{poids}.
\end{solution}

\begin{solution}{exo:g10:inertia:13}
Poussée~: \omterm{def:g10:universal-gravitation:weight}{poids}, réaction, \omterm{def:g10:inertia:inventory}{poussée} de la
main, petit \omterm{def:g10:inertia:inventory}{frottement}~; la vitesse augmente, donc les
\omterm{def:g10:inertia:force}{forces} ne se \omterm{def:g10:inertia:compensate}{compensent} pas
(\cref{prop:g10:inertia:converse}). Glisse~: \omterm{def:g10:universal-gravitation:weight}{poids},
réaction, petit \omterm{def:g10:inertia:inventory}{frottement}~; la pierre ralentit, donc
encore pas de compensation --- le \omterm{def:g10:inertia:inventory}{frottement} est le reste non
compensé (sur une glace idéale elle glisserait pour toujours). Repos~:
\omterm{def:g10:universal-gravitation:weight}{poids} et réaction seuls, se compensant
(\cref{thm:g10:inertia:principle}).
\end{solution}

\begin{solution}{exo:g10:inertia:14}
\emph{1.} Par le principe d'inertie, garder une vitesse constante
n'exige aucune \omterm{def:g10:inertia:force}{force} du tout --- et dans l'espace
interstellaire presque aucune n'agit. Aristote prédirait que la sonde
s'arrête lorsque le «~moteur~» cesse. \emph{2.} Une année
$\approx \qty{3.156e7}{s}$, donc
$d = 17 \times \num{3.156e7} \approx \qty{5.4e8}{km} \approx
\qty{3.6}{au}$. \emph{3.} La vitesse est un vecteur~: changer sa
direction change la vitesse, ce qui exige une \omterm{def:g10:inertia:force}{force} non
compensée.
\end{solution}

\begin{solution}{exo:g10:inertia:15}
Au repos les trois \omterm{def:g10:inertia:force}{forces} se
\omterm{def:g10:inertia:compensate}{compensent}~:
$\vect F_3 = -(\vect F_1 + \vect F_2)$, donc
$\vect F_3\,(-2, -4)$, intensité $\sqrt{4 + 16} = \sqrt{20} \approx
\qty{4.5}{N}$, pointant à l'opposé de la résultante des deux premières
(bas-gauche sur la grille).
\end{solution}

\begin{solution}{pb:g10:inertia:1}
\textbf{1.} \omterm{def:g10:inertia:uniform}{Mouvement rectiligne uniforme} à
\qty{0.75}{m/s}. \omterm{def:g10:inertia:force}{Forces}~: \omterm{def:g10:universal-gravitation:weight}{poids} et la
\omterm{def:g10:inertia:inventory}{réaction normale} de la bande (aucun
\omterm{def:g10:inertia:inventory}{frottement} nécessaire~: la valise ne glisse pas sur la
bande).
\textbf{2.} Oui --- \omterm{def:g10:inertia:uniform}{mouvement rectiligne uniforme}
(\cref{thm:g10:inertia:principle}). $P = 23 \times 9.81 \approx
\qty{226}{N}$, donc $R = \qty{226}{N}$.
\textbf{3.} Les diagrammes sont identiques. Le repos et le
\omterm{def:g10:inertia:uniform}{mouvement rectiligne uniforme} sont un seul et même état
mécanique~: l'inertie ne les distingue pas.
\textbf{4.} La valise garde ses \qty{0.75}{m/s} tandis que la bande
s'arrête~: elle glisse ou bascule en avant jusqu'à ce que le
\omterm{def:g10:inertia:inventory}{frottement} absorbe son mouvement. Rien ne l'a poussée ---
elle a simplement gardé la vitesse qu'elle avait.
\textbf{5.} Les vitesses dans le même sens s'additionnent~:
$0.75 + 1.2 = \qty{1.95}{m/s}$. Sa vitesse est constante, donc les
\omterm{def:g10:inertia:force}{forces} sur elle se \omterm{def:g10:inertia:compensate}{compensent}.
\textbf{6.} $P = 0.160 \times 9.81 \approx \qty{1.57}{N}$ vers le bas,
$R = \qty{1.57}{N}$ vers le haut, et \emph{aucune}
\omterm{def:g10:inertia:force}{force} horizontale.
\textbf{7.} Le palet garde \qty{8.0}{m/s} sans aucun moteur. Les
surfaces du quotidien cachent le \omterm{def:g10:inertia:inventory}{frottement} --- une
\omterm{def:g10:inertia:force}{force} vers l'arrière non compensée --- donc le mouvement
semble mourir de lui-même.
\textbf{8.} Sa vitesse change, donc les \omterm{def:g10:inertia:force}{forces} ne se
\omterm{def:g10:inertia:compensate}{compensent} plus (\cref{prop:g10:inertia:converse})~; le
coupable est le \omterm{def:g10:inertia:inventory}{frottement} de la glace rugueuse.
\textbf{9.} Non~: la vitesse est un \emph{vecteur}, et sa direction
change, donc les \omterm{def:g10:inertia:force}{forces} ne se \omterm{def:g10:inertia:compensate}{compensent}
pas --- la glace pousse latéralement sur ses lames.
\textbf{10.} (a)--(iii), (b)--(i), (c)--(ii).
\textbf{11.} \omterm{def:g10:universal-gravitation:weight}{Poids} $P = 80 \times 9.81 \approx \qty{785}{N}$,
traînée nulle~: rien ne \omterm{def:g10:inertia:compensate}{compense} le
\omterm{def:g10:universal-gravitation:weight}{poids}, donc la vitesse doit changer --- le sauteur
accélère vers le bas.
\textbf{12.} Tant que la traînée $<$ le \omterm{def:g10:universal-gravitation:weight}{poids} les
\omterm{def:g10:inertia:force}{forces} ne se \omterm{def:g10:inertia:compensate}{compensent} pas encore~: la
vitesse continue de croître, et la traînée croît avec elle.
\textbf{13.} Vitesse stable~: \omterm{def:g10:universal-gravitation:weight}{poids} et traînée se
\omterm{def:g10:inertia:compensate}{compensent}, donc la traînée est \qty{785}{N}. Le
mouvement est rectiligne et uniforme --- un corps qui tombe dans
l'exact état mécanique de la valise sur le tapis~: l'inertie déguisée.
\textbf{14.} La traînée dépasse maintenant le
\omterm{def:g10:universal-gravitation:weight}{poids}~: les \omterm{def:g10:inertia:force}{forces} ne se
\omterm{def:g10:inertia:compensate}{compensent} pas, donc la vitesse change --- le
sauteur, encore en descente, ralentit~; à mesure que la vitesse
baisse, la traînée se replie vers le \omterm{def:g10:universal-gravitation:weight}{poids}.
\textbf{15.} Stable à nouveau~: traînée $= P = \qty{785}{N}$ ---
exactement la même qu'à \qty{50}{m/s}.
\textbf{16.} La traînée limite égale toujours le
\omterm{def:g10:universal-gravitation:weight}{poids}, voile ou non. Le parachute change la
\emph{vitesse} à laquelle la traînée atteint \qty{785}{N}~:
\qty{5}{m/s} au lieu de \qty{50}{m/s}.
\textbf{17.} «~\dots les \omterm{def:g10:inertia:force}{forces} agissant sur le corps se
\omterm{def:g10:inertia:compensate}{compensent} (ou aucune n'agit)~» --- les deux lectures
du~\cref{thm:g10:inertia:principle}. Le repos n'est que
$\vect v = \vect 0$, un \omterm{def:g10:inertia:uniform}{mouvement rectiligne uniforme}
comme un autre.
\textbf{18.} La patinoire~: le palet croise sans moteur en vue.
Galilée pointe la surface de plus en plus lisse --- le
\omterm{def:g10:inertia:inventory}{frottement}, non la nature, est ce qui arrête les choses.
\textbf{19.} Des \omterm{def:g10:inertia:force}{forces} agissent ---
\omterm{def:g10:universal-gravitation:weight}{poids} \qty{785}{N} vers le bas et traînée
\qty{785}{N} vers le haut~; elles se
\emph{\omterm{def:g10:inertia:compensate}{compensent}}, c'est précisément pourquoi la
vitesse est constante.
\textbf{20.} Tapis~: \omterm{def:g10:universal-gravitation:weight}{poids} et
\omterm{def:g10:inertia:inventory}{réaction normale}. Patinoire~: \omterm{def:g10:universal-gravitation:weight}{poids} et
\omterm{def:g10:inertia:inventory}{réaction normale}. Parachute~: \omterm{def:g10:universal-gravitation:weight}{poids} et
traînée. La seule loi~: le principe d'inertie, dans ses lectures
directe et réciproque.
\end{solution}
